\documentclass[11pt,a4paper]{article}
\usepackage{fontspec,amsmath,amssymb,amsthm,microtype}
\setmainfont{lmroman10-regular.otf}[BoldFont=lmroman10-bold.otf,ItalicFont=lmroman10-italic.otf,BoldItalicFont=lmroman10-bolditalic.otf]
\usepackage{xcolor}
\usepackage[most]{tcolorbox}
\usepackage{mathrsfs}
\providecommand{\lt}{<}
\providecommand{\gt}{>}
\setmonofont{DejaVuSansMono.ttf}[BoldFont=DejaVuSansMono-Bold.ttf,ItalicFont=DejaVuSansMono-Oblique.ttf,BoldItalicFont=DejaVuSansMono-BoldOblique.ttf,Scale=0.86]
\newfontfamily\leanfallback{STIXTwoMath-Regular.otf}[Scale=0.86]
\definecolor{leanrule}{gray}{0.62}
\definecolor{leanback}{gray}{0.962}
\definecolor{leantitle}{gray}{0.22}
\definecolor{leanmuted}{gray}{0.38}
\newtcolorbox{leanbox}[1]{enhanced,breakable,colback=leanback,colframe=leanrule,
  boxrule=0.4pt,arc=0.6mm,left=2.4mm,right=2.4mm,top=1.8mm,bottom=1.8mm,
  before skip=0.8em,after skip=1em,fonttitle=\small\bfseries,coltitle=leantitle,
  colbacktitle=leanback,titlerule=0.3pt,title={#1}}
\newlength{\leancharwidth}
\newenvironment{leancode}{\par\smallskip\begingroup\ttfamily\footnotesize\raggedright
  \settowidth{\leancharwidth}{M}\setlength{\parskip}{0pt}\setlength{\parindent}{0pt}
  \frenchspacing}{\par\endgroup\smallskip}
\newcommand{\leanline}[2]{\hangindent=\dimexpr#1\leancharwidth+4\leancharwidth\relax
  \hangafter=1\noindent\hspace*{#1\leancharwidth}#2\par}
\newcommand{\leansource}[1]{\par{\footnotesize\color{leanmuted}#1}\par}
\newcommand{\leanlabel}[1]{\par\smallskip\noindent{\small\bfseries #1}\par}
\usepackage[margin=25mm]{geometry}
\usepackage[unicode,hidelinks]{hyperref}
\hypersetup{pdfauthor={LangevinLean Project}}
\setlength{\emergencystretch}{3em}
\setlength{\parskip}{0.4em}
\setlength{\parindent}{0pt}
\allowdisplaybreaks[2]
\raggedbottom
\title{\texorpdfstring{Independence of the entire Brownian future and past}{Independence of the entire Brownian future and past}}
\author{LangevinLean Project}
\date{}
\begin{document}\maketitle
Draft mathematical exposition; these notes may contain mistakes.


At a deterministic cut time \(a\ge0\), all future coordinate increments \((B^i_{a+t}-B^i_a)_{i,t}\) are independent of all past coordinates \((B^i_{t\wedge a})_{i,t}\). Adjoining the independent initial state preserves this whole-path independence.


\section{\texorpdfstring{Regrouping independent pairs into independent families}{Regrouping independent pairs into independent families}}


\textbf{Claim.} Suppose the rows \((X_i,Y_i)\) are mutually independent and, within each row, \(X_i\) and \(Y_i\) are independent. Then the entire families \((X_i)_i\) and \((Y_i)_i\) are independent.



\textbf{Proof.} Regard each row as a two-coordinate independent family. Independence of the rows gives independence of the family indexed by the product of the row set and the two-point coordinate set. The two collections occupy disjoint index sets, so their generated sigma-algebras are independent. This argument does not require countability of the row set.

\begin{leanbox}{Lean: \texttt{iIndep\_\allowbreak{}bool\_\allowbreak{}of\_\allowbreak{}indep}, \texttt{indepFun\_\allowbreak{}families\_\allowbreak{}of\_\allowbreak{}independent\_\allowbreak{}pairs}}
\leansource{\texttt{iIndep\_\allowbreak{}bool\_\allowbreak{}of\_\allowbreak{}indep} \quad WassersteinMixing/\allowbreak{}ActualMarkovNoise.lean, lines 20–36}
\begin{leancode}
\leanline{0}{private\ theorem\ iIndep\_bool\_of\_indep\ \{Ω\ :\ Type*\}\ [MeasurableSpace\ Ω]}
\leanline{4}{(Q\ :\ Measure\ Ω)\ [IsProbabilityMeasure\ Q]\ (m\ :\ Bool\ →\ MeasurableSpace\ Ω)}
\leanline{4}{(hi\ :\ Indep\ (m\ false)\ (m\ true)\ Q)\ :\ iIndep\ m\ Q}
\end{leancode}
\leansource{\texttt{indepFun\_\allowbreak{}families\_\allowbreak{}of\_\allowbreak{}independent\_\allowbreak{}pairs} \quad WassersteinMixing/\allowbreak{}ActualMarkovNoise.lean, lines 41–83}
\begin{leancode}
\leanline{0}{theorem\ indepFun\_families\_of\_independent\_pairs\ \{Ω\ I\ A\ :\ Type*\}}
\leanline{4}{[MeasurableSpace\ Ω]\ [MeasurableSpace\ A]}
\leanline{4}{(Q\ :\ Measure\ Ω)\ [IsProbabilityMeasure\ Q]\ (X\ Y\ :\ I\ →\ Ω\ →\ A)}
\leanline{4}{(hX\ :\ ∀\ i,\ Measurable\ (X\ i))\ (hY\ :\ ∀\ i,\ Measurable\ (Y\ i))}
\leanline{4}{(hrows\ :\ iIndepFun\ (fun\ i\ ω\ ↦\ (X\ i\ ω,\ Y\ i\ ω))\ Q)}
\leanline{4}{(hpair\ :\ ∀\ i,\ IndepFun\ (X\ i)\ (Y\ i)\ Q)\ :}
\leanline{4}{IndepFun\ (fun\ ω\ i\ ↦\ X\ i\ ω)\ (fun\ ω\ i\ ↦\ Y\ i\ ω)\ Q}
\end{leancode}
\leanlabel{How the Lean proof follows the argument}
The Boolean helper expands all finite subsets of a two-point index. \texttt{indep\allowbreak{}Fun\_\allowbreak{}families\_\allowbreak{}of\_\allowbreak{}independent\_\allowbreak{}pairs} combines row and within-row independence by uncurry, then applies independence of sigma-algebras generated by disjoint index subsets.
\end{leanbox}

\section{\texorpdfstring{Whole-path independence on the actual product source}{Whole-path independence on the actual product source}}


\textbf{Claim.} The actual Brownian future at \(a\) is independent of the initial state together with the entire Brownian past.



\textbf{Proof.} For each coordinate, independent Brownian increments give independence of the shifted future path and the path restricted to \([0,a]\). Compose the latter with time clipping. Coordinate paths are mutually independent, so the preceding regrouping result gives vector whole-path independence. Finally take the product with the independent initialization: integrating rectangle events over the initial coordinate leaves the same future marginal and factorizes their masses.

\begin{leanbox}{Lean: \texttt{canonicalBrownian\_\allowbreak{}future\_\allowbreak{}past\_\allowbreak{}independent}, \texttt{indepFun\_\allowbreak{}future\_\allowbreak{}initial\_\allowbreak{}past}, \texttt{actualBrownian\_\allowbreak{}raw\_\allowbreak{}future\_\allowbreak{}initial\_\allowbreak{}past\_\allowbreak{}independent}}
\leansource{\texttt{canonicalBrownian\_\allowbreak{}future\_\allowbreak{}past\_\allowbreak{}independent} \quad WassersteinMixing/\allowbreak{}ActualMarkovNoise.lean, lines 87–115}
\begin{leancode}
\leanline{0}{theorem\ canonicalBrownian\_future\_past\_independent\ (I\ :\ Type*)\ (a\ :\ ℝ≥0)\ :}
\leanline{4}{IndepFun}
\leanline{6}{(fun\ (ω\ :\ ULA.CanonicalBrownian.VectorSample\ I)\ i\ t\ ↦}
\leanline{8}{ULA.CanonicalBrownian.coordinate\ i\ (a\ +\ t)\ ω\ -\ ULA.CanonicalBrownian.coordinate\ i\ a\ ω)}
\leanline{6}{(fun\ (ω\ :\ ULA.CanonicalBrownian.VectorSample\ I)\ i\ t\ ↦}
\leanline{8}{ULA.CanonicalBrownian.coordinate\ i\ (min\ t\ a)\ ω)}
\leanline{6}{(ULA.CanonicalBrownian.vectorProbability\ I)}
\end{leancode}
\leansource{\texttt{indepFun\_\allowbreak{}future\_\allowbreak{}initial\_\allowbreak{}past} \quad WassersteinMixing/\allowbreak{}ActualMarkovNoise.lean, lines 119–139}
\begin{leancode}
\leanline{0}{theorem\ indepFun\_future\_initial\_past\ \{Ω₀\ Ω\ A\ B\ :\ Type*\}}
\leanline{4}{[MeasurableSpace\ Ω₀]\ [MeasurableSpace\ Ω]\ [MeasurableSpace\ A]\ [MeasurableSpace\ B]}
\leanline{4}{(μ\ :\ Measure\ Ω₀)\ (ν\ :\ Measure\ Ω)\ [IsProbabilityMeasure\ μ]\ [IsProbabilityMeasure\ ν]}
\leanline{4}{(X\ :\ Ω\ →\ A)\ (Y\ :\ Ω\ →\ B)\ (hX\ :\ Measurable\ X)\ (hY\ :\ Measurable\ Y)}
\leanline{4}{(hi\ :\ IndepFun\ X\ Y\ ν)\ :}
\leanline{4}{IndepFun\ (fun\ ω\ :\ Ω₀\ ×\ Ω\ ↦\ X\ ω.2)\ (fun\ ω\ :\ Ω₀\ ×\ Ω\ ↦\ (ω.1,\ Y\ ω.2))\ (μ.prod\ ν)}
\end{leancode}
\leansource{\texttt{actualBrownian\_\allowbreak{}raw\_\allowbreak{}future\_\allowbreak{}initial\_\allowbreak{}past\_\allowbreak{}independent} \quad WassersteinMixing/\allowbreak{}ActualMarkovNoise.lean, lines 142–159}
\begin{leancode}
\leanline{0}{theorem\ actualBrownian\_raw\_future\_initial\_past\_independent}
\leanline{4}{(P\ :\ Parameters)\ (V\ :\ E\ P.dimension\ →\ ℝ)\ (hV\ :\ PotentialAssumptions\ P\ V)\ (a\ :\ ℝ≥0)\ :}
\leanline{4}{IndepFun}
\leanline{6}{(fun\ (ω\ :\ ActualSpace\ P)\ i\ t\ ↦}
\leanline{8}{CanonicalSemigroup.brownian\ P\ i\ (a\ +\ t)\ ω\ -\ CanonicalSemigroup.brownian\ P\ i\ a\ ω)}
\leanline{6}{(fun\ (ω\ :\ ActualSpace\ P)\ ↦\ (ω.1,\ fun\ i\ t\ ↦\ CanonicalSemigroup.brownian\ P\ i\ (min\ t\ a)\ ω))}
\leanline{6}{(actualProbability\ P\ V)}
\end{leancode}
\leanlabel{How the Lean proof follows the argument}
\texttt{canonical\allowbreak{}Brownian\_\allowbreak{}future\_\allowbreak{}past\_\allowbreak{}independent} combines the actual coordinate pre-Brownian shift theorem with coordinate-path independence. \texttt{indep\allowbreak{}Fun\_\allowbreak{}future\_\allowbreak{}initial\_\allowbreak{}past} proves the product extension directly by product integration. The final theorem specializes these facts to \texttt{actual\allowbreak{}Probability}.
\end{leanbox}


Earlier results used here: \cite[Section~3.22.1]{LangevinLean2026ULASC}; \cite[Section~3.8.11]{LangevinLean2026ULASC}; \cite[Section~3.8.16]{LangevinLean2026ULASC}; \cite[Section~3.8.17]{LangevinLean2026ULASC}; \cite[Section~3.8.19]{LangevinLean2026ULASC}; \cite[Section~3.8.20]{LangevinLean2026ULASC}; \cite[Section~3.8.21]{LangevinLean2026ULASC}.

\begin{thebibliography}{ULA-LSI}
\bibitem[ULA-SC]{LangevinLean2026ULASC}
LangevinLean Project. \emph{Unadjusted Langevin under Strong Convexity}. LL-01, 2026. Draft report, content revision d34ead06987c. \url{https://odakyildiz.github.io/LangevinLean/docs/pdf/UnadjustedLangevin/FullProof.pdf}.
\end{thebibliography}

\end{document}
